\section{An Executed Continuous-Diffusion Certificate}
\label{sec:r19-continuous-benchmark}
A continuous verification theorem becomes informative only after a candidate, its obstacle inequalities, and a decision budget have been supplied. We give an explicit diffusion benchmark in which these objects can all be verified. This benchmark establishes a continuous-model result in its own right. It neither changes the stored settlement economy nor asserts that its numerical thresholds approximate this benchmark. In particular, it does not discharge the separate constructor error account for the CRRA settlement application.

Let $T=1$, $X_0=0$, and
\begin{equation}
 dX_t=a_t\,dt+dB_t,
 \qquad
 J(a,\tau;F)=\mathbb E\left[-\frac{X_\tau^2}{2}-\frac12\int_0^\tau a_t^2\,dt-F\mathbf1_{\{\tau<1\}}\right].
 \label{eq:r19-lq-economy}
\end{equation}
The agent chooses a stopping time $0\leq\tau\leq1$ and a progressively measurable control with $\mathbb E\int_0^1a_t^2dt<\infty$. Regime 1 permits this control; regime 0 imposes $a=0$. Elective surrender is possible immediately, while forced termination at date 1 carries no fee. The fee does not change the state process. Service is $A=\mathbb E\tau$ and surrender incidence is $H=\Pr(\tau<1)$. The two regimes have the same Brownian volatility, terminal settlement, and fee incidence.

\begin{proposition}[Continuous barriers and exact initial thresholds]
\label{prop:r19-continuous}
The no-surrender values are
\begin{align}
 v_1(t,x)&=-\frac{x^2}{2(2-t)}-\frac12\log(2-t),&
 a^*(t,x)&=-\frac{x}{2-t},\\
 v_0(t,x)&=-\frac{x^2}{2}-\frac{1-t}{2}.&&
 \label{eq:r19-continuous-values}
\end{align}
The uniform and initial thresholds coincide and equal
\begin{equation}
 F^{\rm init}_1=F^{\rm unif}_1=\tfrac12\log2,
 \qquad
 F^{\rm init}_0=F^{\rm unif}_0=\tfrac12.
 \label{eq:r19-continuous-thresholds}
\end{equation}
For fees strictly above the corresponding threshold, every exact optimum operates until date 1. Every global $\eta$-response satisfies $H\leq\eta/(F-F^{\rm init}_r)$ and $A\geq1-H$.
\end{proposition}
\begin{proof}
Direct differentiation gives
\[
 (v_1)_t+\tfrac12(v_1)_{xx}+\sup_a\{a(v_1)_x-a^2/2\}=0,
 \qquad
 (v_0)_t+\tfrac12(v_0)_{xx}=0,
\]
with terminal value $-x^2/2$. Completion of the square identifies $a^*=(v_1)_x$. This linear feedback is admissible. Under any admissible control, $\mathbb E\sup_{t\leq1}X_t^2<\infty$ follows from the Brownian maximal inequality and Cauchy--Schwarz for the drift. Therefore It\^o's formula can be applied up to bounded stopping times by localization; its stochastic integral is square integrable and the quadratic terminal terms are integrable. The nonpositive completed-square residual gives the value upper bound, and the feedback with $\tau=1$ attains it. No compact-domain boundary approximation is used in this argument.

Write $G(x)=-x^2/2$ and $c_1(t)=\tfrac12\log(2-t)$, $c_0(t)=(1-t)/2$. For regime 1,
\[
 v_1(t,x)-(G(x)-F)=\frac{1-t}{2(2-t)}x^2+F-c_1(t)\geq F-c_1(0).
\]
In regime 0 the analogous expression is $F-c_0(t)\geq F-c_0(0)$. Thus at $F\geq c_r(0)$ the continuous candidate dominates every stopping obstacle and is the value with surrender. At $F<c_r(0)$, stopping immediately at $(0,0)$ yields $-F>-c_r(0)$, so the initial no-surrender value is not implemented. This proves both threshold equalities. For $F>c_r(0)$, stopping before 1 incurs an obstacle gap at least $F-c_r(0)$. The stopped verification identity therefore bounds the total private loss below by $(F-c_r(0))\Pr(\tau<1)$, in addition to any control loss. The global loss budget proves the bound on $H$, and $\tau\geq\mathbf1_{\{\tau=1\}}$ proves the bound on $A$.
\end{proof}

\subsection{Participation and a decision-separating budget}
Consider two executable offers, one in each regime, with $F_1=9/25$ and $F_0=51/100$. The agent's outside value is zero. Capacity equals the fee, costs $C(F)=F^2/50$, and is borne by the agent in a separate numeraire. The purchaser pays a nonnegative signing transfer and values expected service at one per unit. The operating state and control problem are unaffected by these transfers. Under exact behavior the minimum transfer is $q_r=c_r(0)+C(F_r)$ and the purchaser obtains
\begin{equation}
 \Pi_1=1-\tfrac12\log2-\frac{(9/25)^2}{50},\qquad
 \Pi_0=1-\tfrac12-\frac{(51/100)^2}{50}.
 \label{eq:r19-continuous-surplus}
\end{equation}
For global $\eta$ behavior, the offer $q_r=c_r(0)+C(F_r)+\eta$ guarantees weak participation, and
\begin{equation}
 \underline\Pi_r(\eta)
 =1-\frac{\eta}{F_r-c_r(0)}-c_r(0)-C(F_r)-\eta
 \label{eq:r19-continuous-budget}
\end{equation}
is a valid purchaser lower bound. The exact return $\Pi_0$ is also an upper bound on a competing regime-0 offer with this same fee and its least transfer, even under favorable service. These comparisons are between the two stated offers, not a claim of global optimality over an unspecified contract space.

\input{revisions/2026-09-20-r19/paper/table_continuous_benchmark}

All numbers in this table are outward enclosures obtained with rational arithmetic. Specifically,
\[
 \log2=2\sum_{j=0}^{m-1}\frac{(1/3)^{2j+1}}{2j+1}+R_m,
 \qquad
 0<R_m\leq\frac{2(1/3)^{2m+1}}{(2m+1)(1-1/9)}.
\]
The checker uses this positive tail bound, not floating-point agreement with a library logarithm. With $m=40$ it verifies strict obstacle slack, both participation inequalities, and $\underline\Pi_1(10^{-6})>\Pi_0$. The PDE residual identities are checked symbolically and proved above; the numerical checker verifies the constants and economic margin. There is no time mesh, state grid, action discretization, terminal approximation, or kernel estimation in this particular certificate. Its transfer-error entries are zero because the displayed analytic functions solve the displayed diffusion, not because finite arrays have been relabeled as continuous.

The benchmark complements, rather than replaces, the original settlement application. It establishes a complete continuous theorem--candidate--verification--decision chain and shows how adjustment can reduce the initial enforcement requirement. Transferring the settlement application's specific magnitudes to its antecedent diffusion still requires bounds for its own constructor, interpolants, controls, and response features. That distinct task is recorded in the model-to-decision account rather than concealed by this exact example.
